% Aqui começa o capítulo do Cronograma
\chapter{Cronograma}\label{chp:cronograma}

Este capítulo apresenta o cronograma estimado para o desenvolvimento completo desta pesquisa, detalhando as etapas já concluídas durante o ano de 2024 e as atividades planejadas para o ano de 2025, culminando na defesa da dissertação de mestrado.

A \Tabela{tab:cronograma_2024} ilustra as principais fases do projeto realizadas ao longo de 2024, desde a pesquisa bibliográfica inicial até a redação preliminar desta dissertação e o desenvolvimento funcional do software proposto.

% Você precisará criar a Tabela do Cronograma 2024 aqui.
% O ideal é usar um pacote como 'ganttchart' ou similar se quiser um gráfico de Gantt,
% ou uma tabela simples se for apenas listagem de atividades por período.
% Dada a imagem que você forneceu, uma tabela simples pode ser mais direta no LaTeX padrão.

\begin{table}[!htp]
\centering
\caption[Cronograma das atividades realizadas em 2024.]{Principais etapas de desenvolvimento da pesquisa concluídas no ano de 2024.}
\label{tab:cronograma_2024}
\begin{tabular}{lp{10cm}} % 'l' para a primeira coluna (mês/período), 'p{width}' para descrição
\toprule
\textbf{Período} & \textbf{Atividade} \\ \midrule
Jan/24 - Mar/24 & Pesquisa Bibliográfica Inicial e Aprofundada \\
Fev/24 - Abr/24 & Aquisição das Imagens e Definição do Processo de Anotação com os Médicos Patologistas \\
Mar/24 - Jun/24 & Etapa de Pré-processamento das Imagens (geração de \textit{patches}, máscaras, normalização) \\
Jun/24 - Jul/24 & Etapa de Limpeza e Organização dos Dados Gerados \\
Jul/24 - Out/24 & Etapas de Treinamento Exploratório dos Modelos de DL \\
Set/24 - Nov/24 & Desenvolvimento das Funcionalidades do Software (aplicação de inferência) \\
Out/24 - Dez/24 & Redação Preliminar da Dissertação para Qualificação \\ \bottomrule
\end{tabular}
\end{table}

Para o ano de 2025, o planejamento das atividades subsequentes da pesquisa, visando a conclusão do mestrado, está sumarizado na \Tabela{tab:cronograma_2025_revisado}. As etapas concentram-se no treinamento final do \textit{ensemble}, na condução dos estudos ablativos e de transferibilidade, na produção científica e na elaboração e defesa da dissertação.

\begin{table}[!htb] % Alterei para htb para maior flexibilidade, !htp também é bom
\centering
\caption[Cronograma das atividades planejadas para 2025.]{Planejamento das principais etapas da pesquisa para o ano de 2025.}
\label{tab:cronograma_2025_revisado} % Novo label para a tabela revisada
\begin{tabular}{l p{0.7\textwidth}} % Ajustei a largura da segunda coluna
\toprule
\textbf{Período} & \textbf{Atividade Principal} \\ 
\midrule
Jan/25 - Mar/25 & Treinamento Final e Validação Robusta do Ensemble de Modelos \\
Abr/25 - Mai/25 & Elaboração e Submissão do Relatório de Qualificação \\
Jun/25 - Jul/25 & Condução de Experimentos Ablativos e Testes de Transferibilidade do \textit{framework} \\
Ago/25 - Set/25 & Redação e Submissão de Artigo Científico (Periódico/Conferência) \\
Out/25 - Nov/25 & Redação Final da Dissertação de Mestrado \\
Dez/25 - Mai/2026         & Defesa da Dissertação de Mestrado \\ 
\bottomrule
\end{tabular}
\end{table}

Este cronograma revisado apresenta as macrotarefas planejadas. Cada uma dessas etapas envolverá atividades mais granulares, como pesquisa bibliográfica complementar contínua, análise detalhada de resultados e preparação para as apresentações formais. O acompanhamento contínuo com a orientação será fundamental para garantir o cumprimento das metas dentro dos prazos estabelecidos.